\section{Successor interface: bounded meta-discovery and scoped negative knowledge}
\label{sec:p1-successor-discovery-failure}

Paper I's separation of object knowledge $K_t$, search-universe/formulation state $W_t$, and governed method state $M_t$ also supplies a boundary for later meta-discovery studies. A surprising trace, recurring failure, semantically distant analogy, or representation proposal is not evidence that a new scientific mechanic exists. It is a typed research proposal that must survive bounded comparison against already available parent mechanisms before any structural claim is considered.

The merged recursive-atom programme now makes this boundary operational. A bounded epistemic-atom study freezes the candidate relation, resource envelope, positive and no-atom controls, strongest parent, protected fields, evaluator identity, interaction hypotheses, and recursion stop rule before outcome-bearing evaluation. The first three broad labels did not survive as indivisible atoms: the tension/opportunity pilot closed as \texttt{MULTIPLE\_TENSION\_ATOMS\_REQUIRED}, the thought-experiment pilot as \texttt{MULTIPLE\_IMAGINATION\_ATOMS\_REQUIRED}, and the concept/primitive pilot as \texttt{MULTIPLE\_CONCEPT\_ATOMS\_REQUIRED}. The common scientific result is therefore the study calculus itself, \texttt{RECURSIVE\_ATOM\_STUDY\_CALCULUS\_FROZEN}, not a universal atom basis or a general creativity mechanism.

The same programme sharpens negative experience. A failure episode records what occurred; reusable failure knowledge is a stronger, scoped object stating what may be excluded, under which responsibility and context conditions, and which changes reopen the question. In the bounded prospective study, that distinction was useful, but the incremental mechanism claim did not survive the strongest parents: the terminal is \texttt{FAILURE\_MEMORY\_USEFUL\_BUT\_STANDARD\_METHODS\_SUFFICIENT}. Paper I therefore uses scoped failure applicability and reopening as governance semantics without claiming a new generic failure-learning algorithm.

Finally, the protected zero-error regime-invention study provides a direct strongest-parent boundary for high-level representation change. Across disjoint 42-case primary and replication splits, ORION-JUMP and the verified representation-regime revision parent made identical decisions and achieved identical protected metrics, with zero false jumps on the no-Jump controls. The registered terminal is \texttt{REPRESENTATION\_INVENTION\_NO\_INCREMENTAL\_VALUE}. This does not show that representation change is unnecessary; it shows that the ORION-specific Jump composition adds no bounded value over the registered strong regime-revision parent on that frozen benchmark.

These successor results reinforce rather than broaden Paper I's authority boundary. Proposal records, atom studies and reopened failures can guide what to test next, but they do not authorize a $K/W/M$ mutation, novelty claim, or open-ended scientific-creativity claim. Paper III owns correspondence and obstruction, Paper IV owns scientific-promotion authority, and Paper V owns adoption and parent-subsumption decisions.
